\documentclass{article}
\usepackage{amsmath,amssymb,amsthm}
\usepackage{algorithm}
\usepackage{algorithmic}
\usepackage{enumitem}
\usepackage{hyperref}
\newtheorem{theorem}{Theorem}
\newtheorem{lemma}{Lemma}
\newtheorem{assumption}{Assumption}
\newtheorem{remark}{Remark}
\newcommand{\prox}{\operatorname{prox}}
\newcommand{\R}{\mathbb{R}}

\begin{document}

\section*{Proximal Gradient Method (PGM)}

\section*{Mathematical Formulation}
Consider the composite optimization problem  

\[
\min_{x\in\R^{d}} \; \Phi(x)\;:=\;f(x)+g(x),
\tag{1}
\]

where  

\begin{itemize}[noitemsep]
  \item $f:\R^{d}\!\to\!\R$ is convex and differentiable with $L$‑Lipschitz continuous gradient,
        \[
        \|\nabla f(x)-\nabla f(y)\|\le L\|x-y\|,\qquad\forall x,y\in\R^{d};
        \tag{2}
        \]
  \item $g:\R^{d}\!\to\!\R\cup\{+\infty\}$ is convex, proper and lower‑semicontinuous (possibly non‑smooth).
\end{itemize}

For a stepsize $\eta>0$, the \emph{proximal operator} of $g$ is  

\[
\prox_{\eta g}(v)\;:=\;\arg\min_{u\in\R^{d}}
\Big\{ g(u)+\frac{1}{2\eta}\|u-v\|^{2}\Big\}.
\tag{3}
\]

The proximal gradient iteration is  

\[
x_{k+1}\;=\;\prox_{\eta g}\!\bigl(x_{k}-\eta\nabla f(x_{k})\bigr),
\qquad k=0,1,2,\dots
\tag{4}
\]

The only hyperparameter is the stepsize $\eta$, which will be restricted to $0<\eta\le 1/L$ for convergence.

\section*{Pseudocode}
\begin{algorithm}[H]
\caption{Proximal Gradient Method}
\begin{algorithmic}[1]
\STATE \textbf{Input:} stepsize $\eta\in(0,1/L]$, initial point $x_{0}\in\R^{d}$, maximum iterations $K$.
\FOR{$k=0$ \textbf{to} $K-1$}
    \STATE Compute the gradient $g_{k}\gets \nabla f(x_{k})$.
    \STATE Take a gradient step $y_{k}\gets x_{k}-\eta g_{k}$.
    \STATE Apply the proximal map $x_{k+1}\gets \prox_{\eta g}(y_{k})$.
\ENDFOR
\STATE \textbf{Output:} $x_{K}$ (or the best iterate w.r.t.\ $\Phi$).
\end{algorithmic}
\end{algorithm}

\section*{Dry Run Example}
We minimise the scalar problem  

\[
f(w)=(w-3)^{2},\qquad g(w)=0,\qquad \eta=0.1,\qquad w_{0}=0 .
\]

Since $g\equiv0$, $\prox_{\eta g}(v)=v$.  
The gradient $\nabla f(w)=2(w-3)$, and $L=2$ (because $f''(w)=2$).  

\[
\begin{array}{c|c|c|c}
k & \nabla f(w_{k}) & w_{k+1}=w_{k}-\eta\nabla f(w_{k}) & \Phi(w_{k+1})\\ \hline
0 & 2(0-3)=-6 & 0-0.1(-6)=0.6 & (0.6-3)^{2}=5.76\\
1 & 2(0.6-3)=-4.8 & 0.6-0.1(-4.8)=1.08 & (1.08-3)^{2}=3.6864\\
2 & 2(1.08-3)=-3.84 & 1.08-0.1(-3.84)=1.464 & (1.464-3)^{2}=2.3660
\end{array}
\]

The objective value decreases at each iteration, illustrating the descent property of (4).

\newpage
\section*{Assumptions}
\begin{assumption}[Convexity and Smoothness of $f$]
\label{ass:f}
$f$ is convex and differentiable with $L$‑Lipschitz gradient (2).
\end{assumption}
\begin{assumption}[Convexity of $g$]
\label{ass:g}
$g$ is convex, proper and lower‑semicontinuous.
\end{assumption}
\begin{assumption}[Stepsize]
\label{ass:eta}
The stepsize satisfies $0<\eta\le 1/L$.
\end{assumption}

\section*{Main Convergence Theorem}
\begin{theorem}[Ergodic $O(1/k)$ rate]\label{thm:ergodic}
Let $\{x_{k}\}_{k\ge0}$ be generated by (4) under Assumptions~\ref{ass:f}–\ref{ass:eta}.  
Define the averaged iterate  

\[
\bar{x}_{K}\;:=\;\frac{1}{K}\sum_{k=0}^{K-1}x_{k}.
\]

Then for any optimal solution $x^{\star}\in\arg\min\Phi$,

\[
\Phi(\bar{x}_{K})-\Phi(x^{\star})
\;\le\;
\frac{\|x_{0}-x^{\star}\|^{2}}{2\eta K}.
\tag{5}
\]

Consequently $\Phi(\bar{x}_{K})\to\Phi(x^{\star})$ at rate $O(1/K)$.
\end{theorem}

\section*{Theoretical Properties}
\begin{itemize}[noitemsep]
  \item \textbf{Stability:} The iteration is a non‑expansive mapping (Lemma~\ref{lem:nonexp}) plus a descent step; thus the sequence $\{x_{k}\}$ remains bounded.
  \item \textbf{Stepsize Sensitivity:} Any $\eta\in(0,1/L]$ guarantees (5). Larger $\eta$ yields a smaller constant $1/(2\eta)$ but may violate the Lipschitz condition.
  \item \textbf{Comparison:} When $g\equiv0$, (4) reduces to the classical gradient method, and (5) coincides with the standard $O(1/k)$ bound. For non‑smooth $g$, the same rate holds without any additional variance or bias terms.
\end{itemize}

\section*{Discussion}
The proof relies solely on convexity, Lipschitz smoothness of $f$, and the \emph{non‑expansiveness} of the proximal operator (Lemma~\ref{lem:nonexp}). No stochasticity or strong convexity is required, which makes the result applicable to a broad class of composite problems such as Lasso, total variation denoising, and support‑vector machine regularisation.

\newpage
\section*{Preliminaries (Definitions and Lemmas)}

\begin{lemma}[Descent Lemma for $L$‑smooth $f$]\label{lem:descent}
For any $x,y\in\R^{d}$,
\[
f(y)\le f(x)+\langle\nabla f(x),y-x\rangle+\frac{L}{2}\|y-x\|^{2}.
\tag{6}
\]
\end{lemma}
\begin{proof}
Standard consequence of $L$‑Lipschitz gradient; see e.g. \cite{nesterov2004introductory}.
\end{proof}

\begin{lemma}[Non‑expansiveness of the proximal operator]\label{lem:nonexp}
For any $u,v\in\R^{d}$ and any $\eta>0$,
\[
\big\|\prox_{\eta g}(u)-\prox_{\eta g}(v)\big\|
\;\le\; \|u-v\|.
\tag{7}
\]
\end{lemma}
\begin{proof}
From the definition (3) and the firm non‑expansiveness of proximal maps (see e.g. \cite{bauschke2011convex}).
\end{proof}

\begin{lemma}[Three‑point inequality]\label{lem:three}
For any $x\in\R^{d}$, $z\in\R^{d}$ and any $\eta>0$,
\[
g\bigl(\prox_{\eta g}(x)\bigr)
\;\le\;
g(z)+\frac{1}{2\eta}\Bigl(\|z-x\|^{2}-\|\prox_{\eta g}(x)-z\|^{2}-\|\prox_{\eta g}(x)-x\|^{2}\Bigr).
\tag{8}
\]
\end{lemma}
\begin{proof}
Directly follows from the optimality condition of (3).  
\end{proof}

\section*{Main Proof (Step‑by‑Step Derivation)}
We prove Theorem~\ref{thm:ergodic}.  
All non‑trivial steps are annotated with \textbf{[Step N]}.

\begin{proof}
Let $x^{\star}$ be any optimal solution of (1).  
Define the gradient step $y_{k}=x_{k}-\eta\nabla f(x_{k})$ and the proximal update $x_{k+1}=\prox_{\eta g}(y_{k})$.

\paragraph{Step 1 (Apply Lemma~\ref{lem:three})}
Using (8) with $x=y_{k}$ and $z=x^{\star}$,
\[
g(x_{k+1})\le g(x^{\star})
+\frac{1}{2\eta}\Bigl(\|x^{\star}-y_{k}\|^{2}
-\|x^{\star}-x_{k+1}\|^{2}
-\|x_{k+1}-y_{k}\|^{2}\Bigr).
\tag{9}
\]
\paragraph{Step 2 (Expand $\|x^{\star}-y_{k}\|^{2}$)}
Since $y_{k}=x_{k}-\eta\nabla f(x_{k})$,
\[
\|x^{\star}-y_{k}\|^{2}
=\|x^{\star}-x_{k}+\eta\nabla f(x_{k})\|^{2}
=\|x^{\star}-x_{k}\|^{2}
+2\eta\langle\nabla f(x_{k}),x^{\star}-x_{k}\rangle
+\eta^{2}\|\nabla f(x_{k})\|^{2}.
\tag{10}
\]

\paragraph{Step 3 (Combine (9) and (10))}
Substituting (10) into (9) yields
\[
\begin{aligned}
g(x_{k+1})\le g(x^{\star})+\frac{1}{2\eta}\Bigl(&\|x^{\star}-x_{k}\|^{2}
+2\eta\langle\nabla f(x_{k}),x^{\star}-x_{k}\rangle
+\eta^{2}\|\nabla f(x_{k})\|^{2}\\
&-\|x^{\star}-x_{k+1}\|^{2}
-\|x_{k+1}-y_{k}\|^{2}\Bigr).
\end{aligned}
\tag{11}
\]

\paragraph{Step 4 (Add $f(x_{k+1})$ and use convexity of $f$)}
Convexity of $f$ gives
\[
f(x^{\star})\ge f(x_{k})+\langle\nabla f(x_{k}),x^{\star}-x_{k}\rangle.
\tag{12}
\]
Adding $f(x_{k+1})$ to both sides of (11) and using (12) we obtain
\[
\begin{aligned}
\Phi(x_{k+1})\;&=\;f(x_{k+1})+g(x_{k+1})\\
&\le f(x_{k})+g(x^{\star})+\frac{1}{2\eta}\Bigl(\|x^{\star}-x_{k}\|^{2}
-\|x^{\star}-x_{k+1}\|^{2}
-\|x_{k+1}-y_{k}\|^{2}\Bigr)\\
&\quad+\frac{\eta}{2}\|\nabla f(x_{k})\|^{2}
+\bigl(f(x_{k+1})-f(x_{k})-\langle\nabla f(x_{k}),x_{k+1}-x_{k}\rangle\bigr).
\end{aligned}
\tag{13}
\]

\paragraph{Step 5 (Apply Descent Lemma (6) to the bracketed term)}
From (6) with $x=x_{k}$, $y=x_{k+1}$,
\[
f(x_{k+1})\le f(x_{k})+\langle\nabla f(x_{k}),x_{k+1}-x_{k}\rangle
+\frac{L}{2}\|x_{k+1}-x_{k}\|^{2}.
\tag{14}
\]
Thus the bracketed term in (13) is bounded by $\frac{L}{2}\|x_{k+1}-x_{k}\|^{2}$.

\paragraph{Step 6 (Relate $\|x_{k+1}-x_{k}\|$ to $\|x_{k+1}-y_{k}\|$)}
Recall $y_{k}=x_{k}-\eta\nabla f(x_{k})$, so
\[
x_{k+1}-x_{k}= (x_{k+1}-y_{k}) + (y_{k}-x_{k})
= (x_{k+1}-y_{k}) -\eta\nabla f(x_{k}).
\tag{15}
\]
Hence,
\[
\|x_{k+1}-x_{k}\|^{2}
= \|x_{k+1}-y_{k}\|^{2}
-2\eta\langle x_{k+1}-y_{k},\nabla f(x_{k})\rangle
+\eta^{2}\|\nabla f(x_{k})\|^{2}.
\tag{16}
\]

\paragraph{Step 7 (Plug (14)–(16) into (13))}
Using (14) to replace the bracketed term and (16) to expand $\|x_{k+1}-x_{k}\|^{2}$, we obtain after cancellation of the mixed inner product terms:

\[
\Phi(x_{k+1})\le \Phi(x^{\star})
+\frac{1}{2\eta}\bigl(\|x^{\star}-x_{k}\|^{2}-\|x^{\star}-x_{k+1}\|^{2}\bigr)
+\Bigl(\frac{\eta L}{2}-\frac{1}{2\eta}\Bigr)\|x_{k+1}-y_{k}\|^{2}.
\tag{17}
\]

\paragraph{Step 8 (Choose stepsize $\eta\le 1/L$)}
Because $\eta\le 1/L$, we have $\frac{\eta L}{2}-\frac{1}{2\eta}\le 0$.  
Thus the last term in (17) is non‑positive and can be dropped:

\[
\Phi(x_{k+1})-\Phi(x^{\star})
\le\frac{1}{2\eta}\Bigl(\|x^{\star}-x_{k}\|^{2}-\|x^{\star}-x_{k+1}\|^{2}\Bigr).
\tag{18}
\]
\textbf{[Step 8]}

\paragraph{Step 9 (Sum over $k=0,\dots,K-1$)}
Summing (18) telescopically yields

\[
\sum_{k=0}^{K-1}\bigl(\Phi(x_{k+1})-\Phi(x^{\star})\bigr)
\le\frac{1}{2\eta}\bigl(\|x^{\star}-x_{0}\|^{2}-\|x^{\star}-x_{K}\|^{2}\bigr)
\le\frac{\|x_{0}-x^{\star}\|^{2}}{2\eta}.
\tag{19}
\]

\paragraph{Step 10 (Use convexity of $\Phi$ for the average iterate)}
Since $\Phi$ is convex,
\[
\Phi(\bar{x}_{K})\;=\;\Phi\!\Bigl(\frac{1}{K}\sum_{k=0}^{K-1}x_{k+1}\Bigr)
\le\frac{1}{K}\sum_{k=0}^{K-1}\Phi(x_{k+1}).
\tag{20}
\]

\paragraph{Step 11 (Combine (19) and (20))}
Dividing (19) by $K$ and invoking (20) gives

\[
\Phi(\bar{x}_{K})-\Phi(x^{\star})
\le\frac{\|x_{0}-x^{\star}\|^{2}}{2\eta K}.
\tag{21}
\]

Equation (21) is exactly the bound (5).  Hence the theorem holds. \qed
\end{proof}

\section*{Rate Derivation (With Constants)}
From (21) we read off the explicit constant $C=\|x_{0}-x^{\star}\|^{2}/(2\eta)$.  
If one chooses the maximal admissible stepsize $\eta=1/L$, the bound becomes

\[
\Phi(\bar{x}_{K})-\Phi(x^{\star})\le\frac{L\|x_{0}-x^{\star}\|^{2}}{2K},
\tag{22}
\]
which matches the classical $O(L/k)$ rate for gradient methods.

\section*{Special Cases}
\begin{enumerate}[noitemsep]
  \item \textbf{Pure gradient descent ($g\equiv0$).}  
        Proximal step reduces to identity, and (4) is the standard gradient iteration. The proof collapses to the classic $O(1/k)$ bound.
  \item \textbf{Indicator function $g=\iota_{C}$ of a closed convex set $C$.}  
        $\prox_{\eta g}$ becomes the Euclidean projection onto $C$. The algorithm is projected gradient descent, and Theorem~\ref{thm:ergodic} recovers the known rate for constrained convex optimisation.
  \item \textbf{Strongly convex $f+g$.}  
        If $\Phi$ is $\mu$‑strongly convex, a refined analysis (omitted for brevity) yields a linear rate $O\big((1-\eta\mu)^{k}\big)$.
\end{enumerate}

\begin{thebibliography}{9}
\bibitem{nesterov2004introductory}
Y.~Nesterov,
\emph{Introductory Lectures on Convex Optimization: A Basic Course},
Springer, 2004.

\bibitem{bauschke2011convex}
H.~H. Bauschke and P.~L. Combettes,
\emph{Convex Analysis and Monotone Operator Theory in Hilbert Spaces},
Springer, 2011.
\end{thebibliography}

\end{document}